\begin{ouverture}[colorimetrie]{3}{Préparer et doser une solution}{Constitution et transformations}{Observer · Modéliser · Expérimenter · Expliquer}
\objectif{Calculer une concentration en masse et une concentration maximale.}
\objectif{Préparer une solution par dissolution ou dilution.}
\objectif{Exploiter une gamme d’étalonnage.}
\prerequis{Proportionnalité, conversions et lecture de graphiques.}
\end{ouverture}
\essentiel
\essSection{1}{Soluté, solvant et concentration}
Dans une solution aqueuse, l'eau joue le rôle de solvant ; une espèce dissoute est un soluté. Le volume final de solution n'est pas le volume d'eau versé avant dissolution.
\begin{definition}[Concentration en masse]$c_m=m_{\mathrm{soluté}}/V_{\mathrm{solution}}$, généralement en \unit{g.L^{-1}}. La masse volumique $\rho=m_{\mathrm{solution}}/V_{\mathrm{solution}}$ utilise la masse totale : c'est une autre grandeur.\end{definition}
À température donnée, la concentration maximale correspond à la saturation. Si on ajoute davantage de soluté que le solvant ne peut en dissoudre, une partie reste solide. La masse ajoutée n'est alors pas intégralement dissoute.
\essSection{2}{Dissoudre et diluer}
Pour une dissolution, calculer $m=c_mV$, peser le soluté, le dissoudre, transférer quantitativement dans la fiole, rincer les récipients et compléter au trait après dissolution et retour à température ambiante. Boucher et homogénéiser.
\schema[corrige]{dissolution}
\begin{loi}[Dilution sans perte de soluté]La masse de soluté dans le prélèvement est conservée : $c_{m,\mathrm{mère}}V_{\mathrm{prélevé}}=c_{m,\mathrm{fille}}V_{\mathrm{fille}}$. Le facteur $F=V_{\mathrm{fille}}/V_{\mathrm{prélevé}}=c_{m,\mathrm{mère}}/c_{m,\mathrm{fille}}$ dépasse 1.\end{loi}
\schema[corrige]{dilution}
\essSection{3}{Un dosage par étalonnage}
On compare une propriété mesurable d'une solution inconnue à celle de solutions de concentrations connues préparées dans les mêmes conditions. Une échelle de teintes permet une estimation visuelle ; une courbe de masse volumique peut permettre une interpolation.
\schema[corrige]{etalonnage}
\begin{attention}[Le domaine de validité]Une droite expérimentale ne justifie pas une extrapolation très éloignée de la gamme. Une solution inconnue trop concentrée doit être diluée avant mesure ; il faut ensuite tenir compte du facteur de dilution.\end{attention}\schema[corrige]{teintes}
    \begin{exemple}[Comparer une teinte]Dans la gamme schématique, X est plus foncée que l'étalon à 4 g/L et plus claire que celui à 6 g/L. On peut proposer $4<c_X<6$ g/L, sans prétendre à une lecture plus précise.\end{exemple}
\essSection{4}{Expérimenter et modéliser}
\begin{experience}[Une dilution au laboratoire]Pour préparer 100,0 mL à partir de 10,0 mL de solution mère, rincer la pipette avec la mère, prélever avec une propipette, transférer dans une fiole jaugée de 100,0 mL, ajouter de l'eau distillée puis ajuster le bas du ménisque au trait à hauteur des yeux. Boucher et retourner plusieurs fois. Ne jamais pipeter à la bouche.\end{experience}
\begin{savoirfaire}[Choisir la verrerie]La pipette jaugée fixe un volume prélevé ; la fiole jaugée fixe un volume final. Le bécher sert à manipuler et dissoudre ; ses graduations ne remplacent pas une verrerie jaugée pour une préparation précise.\end{savoirfaire}
\exercices
\rubrique{Prendre confiance}
\begin{blocLivre}{ex-3-1}\begin{exercice}[1]{Nommer les constituants}
On dissout du glucose dans de l’eau pour obtenir une solution limpide.
\begin{questions}
\item Nommer soluté et solvant.
\item La disparition du solide est-elle une disparition du glucose ?
\end{questions}
\end{exercice}\end{blocLivre}
\begin{blocLivre}{ex-3-2}\begin{exercice}[1]{Concentration directe}
Une solution contient 3,60 g de soluté dans un volume final de 200 mL.
\begin{questions}
\item Convertir le volume en litres.
\item Calculer $c_m$.
\end{questions}
\end{exercice}\end{blocLivre}
\begin{blocLivre}{ex-3-3}\begin{exercice}[1]{Préparer par dissolution}
On souhaite 250,0 mL d’une solution à \qty{8,00}{g.L^{-1}}.
\schema{dissolution}
\begin{questions}
\item Calculer la masse à peser.
\item Décrire les opérations après la pesée.
\end{questions}
\end{exercice}\end{blocLivre}
\begin{blocLivre}{ex-3-4}\begin{exercice}[1]{Une dilution par dix}
On prélève 10,0 mL d’une mère à \qty{40,0}{g.L^{-1}}, puis on complète à 100,0 mL.
\schema{dilution}
\begin{questions}
\item Calculer le facteur de dilution.
\item Déterminer la concentration fille.
\end{questions}
\end{exercice}\end{blocLivre}
\rubrique{S’entraîner}
\begin{blocLivre}{ex-3-5}\begin{exercice}[2]{Saturation}
À la température étudiée, la concentration maximale est \qty{200}{g.L^{-1}}. On prépare 50,0 mL de solution saturée.
\begin{questions}
\item Quelle masse maximale peut être dissoute dans ce volume final ?
\item Que se passe-t-il si 12,0 g ont été introduits pour un volume liquide final fixé à 50,0 mL ?
\end{questions}
\end{exercice}\end{blocLivre}
\begin{blocLivre}{ex-3-6}\begin{exercice}[2]{Choisir le prélèvement}
On prépare 100,0 mL à \qty{3,00}{g.L^{-1}} à partir d’une mère à \qty{15,0}{g.L^{-1}}.
\begin{questions}
\item Calculer le volume à prélever.
\item Choisir la verrerie adaptée.
\end{questions}
\end{exercice}\end{blocLivre}
\begin{blocLivre}{ex-3-7}\begin{exercice}[2]{Lire l’étalonnage}
La masse volumique de la solution inconnue est \qty{1,030}{g.mL^{-1}}. Utiliser la droite fournie.
\schema{etalonnage}
\begin{questions}
\item Lire sa concentration en masse.
\item Cette estimation est-elle une interpolation ?
\end{questions}
\end{exercice}\end{blocLivre}
\begin{blocLivre}{ex-3-8}\begin{exercice}[2]{Deux masses différentes}
100,0 mL d’une solution contiennent 5,00 g de sel. La masse totale de solution est 103,0 g.
\begin{questions}
\item Calculer $c_m$.
\item Calculer $\rho$ et expliquer la différence.
\end{questions}
\end{exercice}\end{blocLivre}
\rubrique{Raisonner et approfondir}
\begin{blocLivre}{ex-3-9}\begin{exercice}[3]{Deux dilutions successives}
Une solution à \qty{80,0}{g.L^{-1}} est diluée par 4 puis par 5.
\begin{questions}
\item Déterminer la concentration finale.
\item Quel est le facteur global ?
\end{questions}
\end{exercice}\end{blocLivre}
\begin{blocLivre}{ex-3-10}\begin{exercice}[3]{Un trait dépassé}
Pour préparer 100,0 mL d’une solution par dissolution, un élève dépasse le trait de jauge.
\begin{questions}
\item La concentration obtenue est-elle trop grande ou trop petite ?
\item Retirer un peu de solution suffit-il à rétablir la concentration ?
\end{questions}
\end{exercice}\end{blocLivre}
\begin{blocLivre}{ex-3-11}\begin{exercice}[3]{Mélanger deux solutions}
On mélange 100 mL à \qty{10,0}{g.L^{-1}} et 200 mL à \qty{4,00}{g.L^{-1}} du même soluté. Les volumes sont supposés additifs.
\begin{questions}
\item Déterminer la masse totale de soluté.
\item Calculer la concentration finale.
\end{questions}
\end{exercice}\end{blocLivre}
\begin{blocLivre}{ex-3-12}\begin{exercice}[3]{Une gamme insuffisante}
La gamme d’étalonnage couvre 0 à \qty{100}{g.L^{-1}}. Une solution inconnue est plus concentrée que le dernier étalon. Après dilution par 5, on mesure \qty{36,0}{g.L^{-1}}.
\begin{questions}
\item Pourquoi la dilution était-elle utile ?
\item Calculer la concentration initiale.
\end{questions}
\end{exercice}\end{blocLivre}

